\chapter{关于本模板}

本模板根据浙江大学研究生院编写的《浙江大学研究生学位论文编写规则》~\cite{zjugradthesisrules}，
在原有的 zjuthesis 模板~\cite{zjuthesis}基础上开发而来。

本模板的本科生版本\cite{zjuthesisrules}得到了浙江大学本科生院老师的支持与审核，
已经在本科生院网上公示。
但当前的研究生版本并未经过研究生院老师的审核，
同学们使用时要注意对照模板与要求，
切不可盲目使用。

作者本人并未编写过浙江大学研究生毕业论文，
所以不清楚具体要求。
如果有热心同学愿意帮忙，
可以替我联系相关老师，我会配合审核并修改代码。

\section{Overleaf 使用注意事项}

如果你在Overleaf上编译本模板，请注意如下事项：

\begin{itemize}
    \item 删除根目录的 ``.latexmkrc'' 文件，否则编译失败且不报任何错误
    \item 字体有版权所以本模板不能附带字体，请务必手动上传字体文件，并在各个专业模板下手动指定字体。
        具体方法参照 GitHub 主页的说明。
        {\textbf 本模板的 Overleaf 版删除了仿宋字体设置，否则会编译超时}。
    \item 当前（2019年9月2日）的Overleaf使用TexLive 2017进行编译，但一些伪粗体复制乱码的问题需要TexLive 2019版本来解决。
        所以各位同学可以在Overleaf上编写论文，但务必使用本地的TexLive 2019来进行最终编译，以免产生查重相关问题。
        具体说明参照 GitHub 主页。
\end{itemize}

\section{节标题}

如\autoref{fig:a-figure}所示，这是一张由TIKZ\cite{tikz}包画出的图片，
我们也可以用includegraphics来插入现有的jpg等格式的图片，
如\autoref{fig:zju-logo}。

\begin{figure}[ht]
    \centering
    \begin{tikzpicture}[thick,help lines/.style={thin,draw=black!50}]
        \def\A{\textcolor{input}{$A$}} \def\B{\textcolor{input}{$B$}}
        \def\C{\textcolor{output}{$C$}} \def\D{$D$}
        \def\E{$E$}

        \colorlet{input}{blue!80!black} \colorlet{output}{red!70!black}
        \colorlet{triangle}{orange}

        \coordinate [label=left:\A] (A) at ($ (0,0) + .1*(rand,rand) $);
        \coordinate [label=right:\B] (B) at ($ (1.25,0.25) + .1*(rand,rand) $);

        \draw [input] (A) -- (B);

        \node [name path=D,help lines,draw,label=left:\D] (D) at (A) [circle through=(B)] {};
        \node [name path=E,help lines,draw,label=right:\E] (E) at (B) [circle through=(A)] {};

        \path [name intersections={of=D and E,by={[label=above:\C]C}}];

        \draw [output] (A) -- (C) -- (B);

        \foreach \point in {A,B,C}
            \fill [black,opacity=.5] (\point) circle (2pt);

        \begin{pgfonlayer}{background}
            \fill[triangle!80] (A) -- (C) -- (B) -- cycle;
        \end{pgfonlayer}
    \end{tikzpicture}

    \caption{\label{fig:a-figure}这是一张图片}
\end{figure}

\begin{figure}[ht]
    \centering
    \includegraphics[width=.4\linewidth]{logo/zju}
    \caption{\label{fig:zju-logo}浙江大学LOGO}
\end{figure}

\par 如\autoref{tab:sample}所示，这是一张自动调节列宽的表格。

\begin{table}[ht]
    \caption{\label{tab:sample}自动调节列宽的表格}
    \begin{tabularx}{\linewidth}{c|X<{\centering}}
        \hline
        第一列 & 第二列 \\ \hline
        xxx & xxx \\ \hline
        xxx & xxx \\ \hline
        xxx & xxx \\ \hline
    \end{tabularx}
\end{table}

\par 如\autoref{equ:sample}，这是一个公式

\begin{equation}
    \label{equ:sample}
    A=\overbrace{(a+b+c)+\underbrace{i(d+e+f)}_{\text{虚数}}}^{\text{复数}}
\end{equation}


\section{小节标题}

\subsection{小小节标题}